\subsection{Summary and Project Positioning}

Despite significant progress in UAV path planning, a gap remains in effectively coupling
physical mobility with protocol-level optimisation. Existing works either optimise the path
while ignoring the protocol's failure modes (creating the SF Monoculture) or optimise the
protocol using static mathematics that cannot adapt to mobility. Explicitly acknowledging this
gap, Zorbas et al.~\cite{zorbas2025revisiting} conclude that ``reinforcement learning models
could be employed to adapt over time to evolving network conditions.''

Table~\ref{tab:related_work} positions this dissertation at the specific intersection no
prior work occupies: protocol-aware DRL for mobile LoRaWAN gateways with explicit Capture
Effect modelling and multi-objective fairness.

\begin{table}[H]
\centering
\caption{Positioning of this work relative to key related works across five dimensions.
         \checkmark~= addressed; \texttimes~= not addressed.}
\label{tab:related_work}
\begin{tabularx}{\textwidth}{l Z Z Z Z Z}
\toprule
Work & Protocol- & DRL- & LoRaWAN & Capture & Multi- \\
     & Aware & Based & Specific & Effect & Objective \\
\midrule
Zorbas et al.~\cite{zorbas2025revisiting} & \checkmark & \texttimes & \checkmark & \texttimes & \texttimes \\
Xiong et al.~\cite{xiong2023flyinglora}   & \texttimes & \texttimes & \checkmark & \texttimes & \checkmark \\
Challita et al.~\cite{challita2018interference} & \checkmark & \checkmark & \texttimes & \texttimes & \texttimes \\
Nemer et al.~\cite{nemer2022energyefficient} & \texttimes & \checkmark & \texttimes & \texttimes & \checkmark \\
Luong et al.~\cite{luong2019drl} & \texttimes & \checkmark & \texttimes & \texttimes & \texttimes \\
\textbf{This work} & \checkmark & \checkmark & \checkmark & \checkmark & \checkmark \\
\bottomrule
\end{tabularx}
\end{table}

In this report, the focus is placed on this specific intersection: using a DRL agent to
solve the SF Monoculture. Unlike the static optimisation in~\cite{zorbas2025revisiting} or
the energy-centric heuristics in~\cite{xiong2023flyinglora}, the proposed framework treats
the UAV's physical position as a direct control input for the LoRaWAN protocol, allowing the
agent to learn an emergent policy that balances energy efficiency with rigorous protocol
health. The Capture Effect model distinguishes this work from all prior LoRaWAN
simulators that assume mutual packet destruction on any intra-SF collision.
